\begin{Exercice}
On considère les fonctions numériques $f$ et $g$ définies par : \quad
$f(x) = -x^2 + 4x - 3  \quad \text{et} \quad  g(x) = \frac{x-3}{x-2}$
\vspace{0.2cm} \\
On note $(C_f)$ et $(C_g)$ leurs courbes représentatives dans un repère orthonormé $(O;\vec{i};\vec{j})$.

\begin{enumerate}
\item
\begin{enumerate}
    \item Déterminer la nature de $(C_f)$ en précisant ses éléments caractéristiques.
    \item Calculer $f(0)$, $f(1)$ et $f(3)$.
    \item Dresser le tableau de variations de la fonction $f$ sur $\mathbb{R}$.
\end{enumerate}
\item
\begin{enumerate}
    \item Déterminer $D_g$, l'ensemble de définition de la fonction $g$.
    \item Dresser le tableau de variations de la fonction $g$ sur $D_g$.
    \item Déterminer la nature de $(C_g)$ en précisant ses éléments caractéristiques.
\end{enumerate}
\item
\begin{enumerate}
    \item Montrer que pour tout $x \in D_g$, on a :
    $f(x) - g(x) = \frac{(x-3)(-x^2+3x-3)}{x-2}$
    \item Montrer que pour tout réel $x$ : $-x^2+3x-3 < 0 $
    \item En déduire les coordonnées du point d'intersection des courbes $(C_f)$ et $(C_g)$.
\end{enumerate}
\item Tracer les asymptotes, la courbe $(C_g)$ et la courbe $(C_f)$ dans le même repère $(O;\vec{i};\vec{j})$.
\item  Par lecture graphique, résoudre dans $\mathbb{R}$ l'inéquation :
$\frac{x-3}{x-2} \le -x^2 + 4x - 3$
\item  Discuter graphiquement, suivant les valeurs du paramètre réel $m$, le nombre de solutions de l'équation :$$-x^2 + 4x - 3 - m = 0$$
\end{enumerate}
\end{Exercice}
